

%----------------------------------------------------------------------------------------------------------
%	Contents
% ----------------------------------------------------------------------------------------------------------
%

\section{Conclusion}
\subsection{Discussion}
catch the main motivation from the introduction and try to mesh the results from the study into the outlook in regards to HCI. What was initially explored, how did we achieve it and how can we conclude the results.
\subsection{Limitations} 
\begin{itemize}
\item Kinect and its sdk
\begin{itemize}
\item has some flaws regarding infrared interferences
\item tons of noises that needs to be filtered correctly
\end{itemize}
\item Our application
\begin{itemize}
\item Unwrap algorithm is not yet capable to deal with more complex meshes
\item manipulation of camera angles sometimes feel clunky (system dependent)
\end{itemize}
\end{itemize}
\subsection{Future work}
\begin{itemize}
\item Connections and wiring of on-body sensors is a challenging task. The digital tools could provide suggestions or intelligently route the connections of the sensors taking into account the geometric properties of the body.
\item great to prototype sensors for any object shape i.e. on prosthetics
\item can also be used to create individualistic 3d objects i.e. a tailor could digitally design a garment directly on a customer’s body, rather than manually measuring the body dimensions and subsequently producing the design.
\item design with your fingers
\end{itemize}
